\documentclass[11pt]{article}
\usepackage[margin=1in]{geometry}
\usepackage{array}
\usepackage{enumitem}
\usepackage{xcolor}
\usepackage{hyperref}
\usepackage{listings}

\hypersetup{colorlinks=true, linkcolor=blue, urlcolor=blue}
\setlength{\parindent}{0pt}
\setlength{\parskip}{6pt}
\setlist[itemize]{topsep=3pt,itemsep=2pt,leftmargin=*}
\setlist[enumerate]{topsep=3pt,itemsep=4pt,leftmargin=*}

\lstset{
  basicstyle=\ttfamily\small,
  breaklines=true,
  frame=single,
  columns=fullflexible,
  keepspaces=true,
  showstringspaces=false,
  backgroundcolor=\color{gray!6}
}

\definecolor{predictcolor}{RGB}{255,242,200}
\definecolor{discusscolor}{RGB}{214,234,255}
\definecolor{funfactcolor}{RGB}{222,247,222}
\definecolor{debatecolor}{RGB}{255,221,221}

\newcommand{\calloutbox}[3]{%
  \vspace{4pt}\noindent
  \colorbox{#1}{\parbox{0.96\linewidth}{\textbf{#2}}}\\[2pt]
  \noindent\fbox{\begin{minipage}{0.94\linewidth}\vspace{2pt}#3\vspace{2pt}\end{minipage}}
  \vspace{6pt}
}
\newcommand{\predictbox}[1]{\calloutbox{predictcolor}{PREDICTION TIME}{#1}}
\newcommand{\discussbox}[1]{\calloutbox{discusscolor}{HUDDLE UP}{#1}}
\newcommand{\funfact}[1]{\calloutbox{funfactcolor}{FUN FACT}{#1}}
\newcommand{\debatebox}[1]{\calloutbox{debatecolor}{DEBATE CORNER}{#1}}
\newcommand{\claimbox}[1]{\calloutbox{funfactcolor}{CLAIM, EVIDENCE, CONTEXT}{#1}}
\newcommand{\badge}[1]{$[\ ]$ \textbf{#1}}

\newcommand{\writebox}[1]{
  \vspace{3pt}
  \noindent\fbox{
    \begin{minipage}{0.96\linewidth}
      \textbf{#1}\\[12pt]
      \rule{\linewidth}{0.4pt}\\[14pt]
      \rule{\linewidth}{0.4pt}\\[14pt]
      \rule{\linewidth}{0.4pt}\\[14pt]
    \end{minipage}
  }
  \vspace{6pt}
}

\newcommand{\shortwritebox}[1]{
  \vspace{3pt}
  \noindent\fbox{
    \begin{minipage}{0.96\linewidth}
      \textbf{#1}\\[12pt]
      \rule{\linewidth}{0.4pt}\\[14pt]
    \end{minipage}
  }
  \vspace{6pt}
}

\begin{document}

\begin{center}
  {\LARGE MA 213 Week 6: Lab 5}\\[3pt]
  {\large Sampling Distributions and the Central Limit Theorem}
\end{center}

\begin{enumerate}
  \item \textbf{Your name:} \underline{\hspace{0.3\textwidth}} \hfill \textbf{BU ID:} \underline{\hspace{0.25\textwidth}}
  \item \textbf{Partner's name:} \underline{\hspace{0.3\textwidth}} \hfill \textbf{BU ID:} \underline{\hspace{0.25\textwidth}}
\end{enumerate}

\textbf{Big idea:} A sample statistic changes from sample to sample. In this lab, you will use simulation to see the sampling distribution of a sample mean and explain how the Central Limit Theorem makes inference possible.
\textbf{Do not use GenAI tools for this lab.} The point is to practice your own simulation work, statistical reasoning, and interpretation.

\textbf{Files used:} No data file is needed for the main activity. You will generate hypothetical populations in R.

\textbf{Skills practiced:} \texttt{sample()}, \texttt{mean()}, \texttt{sd()}, \texttt{for}, \texttt{function()}, \texttt{rnorm()}, \texttt{rbinom()}, \texttt{rpois()}, \texttt{rexp()}, \texttt{rgamma()}, and \texttt{ggplot()}.

\textbf{Your mission:} Build a population, repeatedly sample from it, graph the sampling distribution of $\hat{\mu}$, and explain how changing $n$ and $K$ changes what you see.

\section*{Icebreaker}

Before opening R: introduce yourself to your partner. Then answer: If you sampled 10 BU students and recorded their commute time, would your sample mean be exactly the same as another group's sample mean? Why?

\shortwritebox{My answer and my partner's answer:}

\section*{Getting Started: Population, Sample, Statistic}

A \textbf{population} is the full group we care about. A \textbf{sample} is a smaller group selected from the population. A \textbf{statistic} is a number computed from a sample, such as the sample mean $\hat{\mu}$.

\discussbox{A sampling distribution is not the distribution of the raw data. It is the distribution of a statistic computed from many repeated samples. Keep that distinction in mind today.}

\section*{Question 1. Build your population}

Choose one population below. Different groups should choose different populations when possible.

\begin{center}
\begin{tabular}{|p{0.25\textwidth}|p{0.32\textwidth}|p{0.33\textwidth}|}
\hline
\textbf{Population} & \textbf{R generator} & \textbf{Parameters}\\
\hline
Normal & \texttt{rnorm()} & $\mu = 2$, $\sigma = 0.447$\\[6pt]
\hline
Bernoulli & \texttt{rbinom()} & $p = 0.3$\\[6pt]
\hline
Poisson & \texttt{rpois()} & $\lambda \approx 0.2$\\[6pt]
\hline
Exponential & \texttt{rexp()} & rate $\approx 2.237$\\[6pt]
\hline
Gamma & \texttt{rgamma()} & shape $=4$, rate $\approx 4.474$\\[6pt]
\hline
\end{tabular}
\end{center}

\begin{lstlisting}[language=R]
library(dplyr)
library(ggplot2)

set.seed(213)
sigma <- 0.447
N <- 100000

# Choose ONE population and comment out the others.
pop <- rnorm(N, mean = 2, sd = sigma)
# pop <- rbinom(N, size = 1, prob = 0.3)
# pop <- rpois(N, lambda = sigma^2)
# pop <- rexp(N, rate = 1 / sigma)
# pop <- rgamma(N, shape = 4, rate = sqrt(4) / sigma)

population_mean <- mean(pop)
population_sd <- sd(pop)

population_mean
population_sd

ggplot(data.frame(X = pop), aes(x = X)) +
  geom_histogram(bins = 40, fill = "steelblue", color = "white") +
  labs(
    x = "Population value",
    y = "Count",
    title = "My hypothetical population"
  )
\end{lstlisting}

\writebox{Which population did you choose? Describe its shape, center, and spread.}

\section*{Question 2. Draw one sample}

Draw one sample of size $n=100$ from your population and compute the sample mean $\hat{\mu}$.

\begin{lstlisting}[language=R]
n <- 100

my_sample <- sample(pop, size = n)
mu_hat <- mean(my_sample)

mu_hat
population_mean
mu_hat - population_mean
\end{lstlisting}

\predictbox{Before comparing with other groups, do you expect your \texttt{mu\_hat} to equal the population mean exactly? Why or why not?}

\writebox{How close was your sample mean to the population mean?}

\section*{Question 3. Build a sampling distribution}

Now repeat the sampling process $K=100$ times. Each repetition gives one sample mean. The vector of sample means approximates the sampling distribution of $\hat{\mu}$.

\begin{lstlisting}[language=R]
K <- 100
n <- 100

mu_hat_vector <- rep(0, K)

for (i in 1:K) {
  my_sample <- sample(pop, size = n)
  mu_hat_vector[i] <- mean(my_sample)
}

head(mu_hat_vector)
mean(mu_hat_vector)
sd(mu_hat_vector)
\end{lstlisting}

\begin{lstlisting}[language=R]
ggplot(data.frame(mu_hat = mu_hat_vector), aes(x = mu_hat)) +
  geom_histogram(bins = 20, fill = "steelblue", color = "white") +
  geom_vline(xintercept = population_mean, color = "red", linetype = "dashed") +
  labs(
    x = "Sample mean",
    y = "Count",
    title = "Sampling distribution of the sample mean",
    subtitle = "Red dashed line shows the population mean"
  )
\end{lstlisting}

\discussbox{Where is the sampling distribution centered? Is it centered near one sample mean, or near the population mean?}

\writebox{Describe the shape, center, and spread of your first sampling distribution.}

\section*{Question 4. Write a reusable simulation function}

Wrap the repeated sampling process into a function. This lets you quickly compare different sample sizes and repetition counts.

\begin{lstlisting}[language=R]
get_mu_hats <- function(pop, n, K) {
  mu_hat_vector <- rep(0, K)

  for (i in 1:K) {
    my_sample <- sample(pop, size = n)
    mu_hat_vector[i] <- mean(my_sample)
  }

  return(mu_hat_vector)
}

result_q4 <- get_mu_hats(pop, n = 100, K = 100)
head(result_q4)
\end{lstlisting}

\writebox{How could you check whether \texttt{result\_q4} behaves like your vector from Question 3?}

\section*{Question 5. What does sample size $n$ change?}

Keep $K=1000$ fixed. Compare the sampling distributions when $n=20$, $n=200$, and $n=1000$.

\begin{lstlisting}[language=R]
mu_hats_20 <- get_mu_hats(pop, n = 20, K = 1000)
mu_hats_200 <- get_mu_hats(pop, n = 200, K = 1000)
mu_hats_1000 <- get_mu_hats(pop, n = 1000, K = 1000)

sampling_n <- data.frame(
  mu_hat = c(mu_hats_20, mu_hats_200, mu_hats_1000),
  sample_size = factor(rep(c(20, 200, 1000), each = 1000))
)

ggplot(sampling_n, aes(x = mu_hat)) +
  geom_histogram(bins = 30, fill = "steelblue", color = "white") +
  facet_wrap(~ sample_size, scales = "free_y") +
  geom_vline(xintercept = population_mean, color = "red", linetype = "dashed") +
  labs(
    x = "Sample mean",
    y = "Count",
    title = "Effect of sample size on the sampling distribution"
  )
\end{lstlisting}

\begin{center}
\begin{tabular}{|p{0.16\textwidth}|p{0.22\textwidth}|p{0.22\textwidth}|p{0.22\textwidth}|}
\hline
\textbf{Sample size} & \textbf{Shape} & \textbf{Center} & \textbf{Spread}\\
\hline
$n=20$ & & & \\[22pt]
\hline
$n=200$ & & & \\[22pt]
\hline
$n=1000$ & & & \\[22pt]
\hline
\end{tabular}
\end{center}

\claimbox{As $n$ increases, the sampling distribution of $\hat{\mu}$ becomes \underline{\hspace{4cm}} and its spread becomes \underline{\hspace{3cm}}.}

\section*{Question 6. What does number of repetitions $K$ change?}

Now keep $n=1000$ fixed and change $K$. This changes how clearly you can see the sampling distribution, but it does not change the theoretical spread of $\hat{\mu}$.

\begin{lstlisting}[language=R]
mu_hats_k20 <- get_mu_hats(pop, n = 1000, K = 20)
mu_hats_k200 <- get_mu_hats(pop, n = 1000, K = 200)
mu_hats_k1000 <- get_mu_hats(pop, n = 1000, K = 1000)

sampling_k <- data.frame(
  mu_hat = c(mu_hats_k20, mu_hats_k200, mu_hats_k1000),
  repetitions = factor(c(rep(20, 20), rep(200, 200), rep(1000, 1000)))
)

ggplot(sampling_k, aes(x = mu_hat)) +
  geom_histogram(bins = 25, fill = "coral", color = "white") +
  facet_wrap(~ repetitions, scales = "free_y") +
  geom_vline(xintercept = population_mean, color = "red", linetype = "dashed") +
  labs(
    x = "Sample mean",
    y = "Count",
    title = "Effect of number of repetitions on the picture"
  )
\end{lstlisting}

\debatebox{Increasing $n$ and increasing $K$ both make the graphs look different. Which one changes the actual sampling variability? Which one mostly changes how clearly we can see the distribution?}

\writebox{Explain the difference between increasing $n$ and increasing $K$.}

\section*{Question 7. Compare standard error and simulated spread}

The standard error of the sample mean is
\[
SE = \frac{\sigma}{\sqrt{n}}.
\]
Compare the theoretical standard error with the empirical standard deviation of your simulated sample means.

\begin{lstlisting}[language=R]
n <- 1000
K <- 1000

mu_hats <- get_mu_hats(pop, n = n, K = K)

SE <- population_sd / sqrt(n)
empirical_sd <- sd(mu_hats)

SE
empirical_sd
abs(SE - empirical_sd)
\end{lstlisting}

\writebox{Are the theoretical standard error and simulated standard deviation close? Explain in context.}

\section*{Optional Challenge}

Use the empirical rule to check whether your sampling distribution is approximately Normal.

\begin{lstlisting}[language=R]
m <- mean(mu_hats)
s <- sd(mu_hats)

rate1 <- mean(abs(mu_hats - m) < 1 * s)
rate2 <- mean(abs(mu_hats - m) < 2 * s)
rate3 <- mean(abs(mu_hats - m) < 3 * s)

rate1
rate2
rate3

abs(c(rate1 - 0.68, rate2 - 0.95, rate3 - 0.997))
\end{lstlisting}

\writebox{Does your sampling distribution approximately follow the 68--95--99.7 rule?}

\section*{Final Reflection}

\begin{enumerate}
  \item What is a sampling distribution in your own words?
  \item What happened to shape, center, and spread as $n$ increased?
  \item What changed when $K$ increased?
  \item Did your original population shape make it easy or hard for the CLT to appear?
\end{enumerate}

\writebox{Write one paragraph summarizing what you learned from the simulation.}

\end{document}
